%!TEX root = ../main.tex

\chapter{Knowledge Graph Retrieval}\label{ch:kgret}

This chapter uses the graph of Chapter~\ref{ch:kg} for conversational question answering.
Every pipeline is evaluated on the dev-500 subset with the generator, prompt and
metrics of Chapter~\ref{ch:rag}, so its numbers are directly comparable to passage
retrieval. A graph retriever has to start from entities: the user turn is first linked
to entities of the graph, and the linked entities then select sections, facts, or both.
Section~\ref{sec:kgret:linking} describes the entity linker and
Section~\ref{sec:kgret:linkeval} evaluates it on its own, before any answer is
generated. Section~\ref{sec:kgret:status} summarises the decisions it implies for the
graph retrievers.


% ===========================================================================
\section{Entity Linking}\label{sec:kgret:linking}

\paragraph{Input.} Only 28\% of the dev-500 questions name an entity of their gold
section (Section~\ref{sec:kg:full}), so the current user turn alone is a weak input.
Three variants are compared: the raw user turn; the standalone query produced by the
rewriter of Section~\ref{sec:rag:rewrite}; and the rewritten query followed by the last
agent turn. The rewrites reuse the prompt and the cache of Chapter~\ref{ch:rag}, so
linking adds no \gls{llm} calls: there are 0.87 calls per question, the rewrite itself,
which is skipped on first turns.

\paragraph{Candidates.} Each entity is represented by its surface forms: its name, the
names of the mentions resolved to it, and their aliases (19{,}067 forms for 14{,}731
entities). Two candidate generators score the entities:
\begin{itemize}
    \item \textbf{dense}: the bge-small model of the dense retriever embeds the query
    (with the BGE query instruction) and every surface form, and an entity scores the
    cosine of its best form;
    \item \textbf{\gls{bm25}} over one short document per entity, made of its surface
    forms without stopwords.
\end{itemize}
The top-100 of each generator are fused with \gls{rrf} ($k = 60$, as in the hybrid
retriever).

\paragraph{Filters.} Two optional filters restrict the candidates.
\begin{itemize}
    \item \textbf{Domain.} The retrieval pipeline of Chapter~\ref{ch:rag} has no domain
    classifier, so the predicted domain is the majority domain of the top-5 sections of
    the best passage retriever (hybrid, rewrite, reranker), and its confidence is the
    share of the top-5 in that domain. When at least 4 of the 5 sections agree, only
    entities seen in that domain are kept. The prediction is correct for 94.8\% of the
    dev-500 questions, is confident for 471 of them, and is correct for 97.9\% of the
    confident ones.
    \item \textbf{Hubs.} Entities with a degree above the 99th percentile of the graph
    (more than 16 relations; 126 entities, for example ``VA health care'' with 99) are
    skipped.
\end{itemize}
Both thresholds were fixed before the evaluation.


% ===========================================================================
\section{Linking Evaluation}\label{sec:kgret:linkeval}

\paragraph{Metric.} Link recall@$k$ is the share of examples in which at least one of
the top-$k$ linked entities has a mention in a gold section. All 500 examples have a
gold section in the grounding annotation, including the 31 unanswerable turns; for
96.6\% of them the gold section has at least one entity, which bounds link recall. The
metric is lenient, since an entity mentioned in many sections reaches the gold section
more easily. It is therefore reported with two controls:
\begin{itemize}
    \item \textbf{fan-out}: the number of distinct sections that the top-10 entities are
    mentioned in;
    \item \textbf{chance}: the recall obtained when each example receives the linked
    entities of another random example of the same domain.
\end{itemize}
Intervals are 95\% Wilson intervals.

\begin{table}[ht]
    \centering
    \caption{Link recall on dev-500 ($n = 500$, \%). Default linker: dense and
    \gls{bm25} fused by \gls{rrf}, domain filter, hub skip. Ablations use the
    rewritten query.}
    \label{tab:kgret:linking}
    \begin{tabular}{lcccrr}
        \toprule
        Configuration & R@1 & R@5 & R@10 & Chance & Fan-out \\
        \midrule
        Raw turn                 & 24.2 & 41.2 & 47.4 & 3.6 & 20 \\
        Rewritten                & 29.6 & 52.2 & 62.0 & 5.1 & 22 \\
        Rewritten + agent turn   & 24.8 & 49.2 & 59.4 & 5.4 & 21 \\
        \midrule
        Dense only               & 27.2 & 49.2 & 59.2 & 5.3 & 21 \\
        \gls{bm25} only          & 32.8 & 50.0 & 59.4 & 5.0 & 20 \\
        No domain filter         & 29.6 & 51.6 & 61.8 & 4.8 & 21 \\
        Gold domain              & 30.6 & 54.2 & 63.8 & 5.2 & 22 \\
        \textbf{No hub skip}     & \textbf{33.2} & \textbf{59.6} & \textbf{70.4} & 7.0 & 29 \\
        Raw + rewritten fused    & 26.0 & 47.6 & 57.8 & 4.7 & 21 \\
        \bottomrule
    \end{tabular}
    \\[2pt]
    {\footnotesize Chance: recall@10 with the linked entities of a random example of the
    same domain (mean of 5 permutations). Fan-out: median number of sections reached by
    the top-10 entities.}
\end{table}

Table~\ref{tab:kgret:linking} summarises the results.
\begin{itemize}
    \item \textbf{Rewrite.} Rewriting the turn is the largest single gain: recall@10
    rises from 47.4\% (CI $[43.1, 51.8]$) to 62.0\% ($[57.7, 66.1]$). Appending the last
    agent turn does not help and lowers recall@1 by five points, because the previous
    answer contributes its own entities.
    \item \textbf{Generators and domain.} Dense and \gls{bm25} candidates are
    complementary: \gls{bm25} is best at rank 1, where exact name overlap matters, and
    their fusion is best at ranks 5 and 10. The domain filter is almost neutral: even
    the gold domain adds less than two points.
    \item \textbf{Hubs.} Skipping hubs is the one clearly harmful choice. Keeping them
    raises recall@10 to 70.4\% ($[66.3, 74.2]$), with 46 more examples linked. The median
    fan-out grows from 22 to 29 sections, but chance recall stays at 7\%, so the extra
    hits are genuine: many questions are about a hub itself, such as ``VA health care''
    or ``Social Security benefits''. This holds most in SSA, which gains 13.6 points.
\end{itemize}

\paragraph{Document switches.} Turns that move to a new document behave differently
(Table~\ref{tab:kgret:splits}). There the raw turn links better than the rewrite
(62.8\% vs 55.3\% at recall@10): a turn that opens a new topic usually names it, while
the rewriter tends to carry the previous topic over. The last agent turn makes this
worse (42.6\%). Fusing the rankings of the raw and the rewritten query recovers the
switch turns (66.0\%) but loses on the much larger set of same-document turns. With 94
switch turns, these differences lie within the confidence intervals.

\begin{table}[ht]
    \centering
    \caption{Link recall@10 by turn type (\%, 95\% Wilson interval).}
    \label{tab:kgret:splits}
    \small
    \begin{tabular}{lccc}
        \toprule
        Configuration & Same document (406) & Doc switch (94) & Unanswerable (31) \\
        \midrule
        Raw turn               & 43.8 {\scriptsize[39.1, 48.7]} & 62.8 {\scriptsize[52.7, 71.9]} & 16.1 {\scriptsize[7.1, 32.6]} \\
        Rewritten              & 63.5 {\scriptsize[58.8, 68.1]} & 55.3 {\scriptsize[45.3, 65.0]} & 67.7 {\scriptsize[50.1, 81.4]} \\
        Rewritten + agent turn & 63.3 {\scriptsize[58.5, 67.8]} & 42.6 {\scriptsize[33.0, 52.6]} & 80.6 {\scriptsize[63.7, 90.8]} \\
        Rewritten, hubs kept   & 72.2 {\scriptsize[67.6, 76.3]} & 62.8 {\scriptsize[52.7, 71.9]} & 80.6 {\scriptsize[63.7, 90.8]} \\
        Raw + rewritten, hubs kept & 64.8 {\scriptsize[60.0, 69.3]} & 66.0 {\scriptsize[55.9, 74.7]} & 58.1 {\scriptsize[40.8, 73.6]} \\
        \bottomrule
    \end{tabular}
\end{table}

\paragraph{Errors.} Of the 148 examples that the best configuration misses at
recall@10, 17 have no entity in their gold section. In most of the others, the linker
finds entities on the right topic that are different nodes from those in the gold
section. For example, it links ``disability benefits'' and ``enough Social Security
credits at this time to qualify'', while the gold section mentions ``monthly disability
benefit'' and ``enough credits''. Long, clause-like entity names and conservative entity
resolution split one concept into several nodes. One hop through entity relations, or
section scoring by entity similarity rather than identity, can recover part of these
cases. The remaining misses are turns without content words (``do not'') and gold
sections that list many items, whose entities are the items rather than the topic the
question names.

\paragraph{Cost.} Linking takes about 15\,ms per question on a CPU (6\,ms query
encoding, 9\,ms scoring). It needs no \gls{llm} calls beyond the rewrite.


% ===========================================================================
\section{Status}\label{sec:kgret:status}

Entity linking reaches the gold section for 70\% of the dev-500 questions with the top-10
entities, about the level of section recall@5 of the best passage retriever (67.4\%).
However, a linked entity is mentioned in a median of 29 sections, so ranking these
sections, rather than linking, will decide how graph retrieval compares with passage
retrieval. The graph retrievers therefore use:
\begin{itemize}
    \item the rewritten query, without the agent turn;
    \item the top-10 entities;
    \item the domain filter;
    \item hubs kept in linking, but down-weighted by the number of sections they are
    mentioned in when sections are scored, and excluded from relation expansion.
\end{itemize}
Three graph retrievers are compared next against the best passage retriever:
\begin{itemize}
    \item graph to sections;
    \item graph facts as context;
    \item a hybrid of graph and passage retrieval.
\end{itemize}
